% Part: normal-modal-logic
% Chapter: completeness
% Section: introduction

\documentclass[../../../include/open-logic-section]{subfiles}

\begin{document}

\olfileid{nml}{com}{int}

\olsection{પરિચય}

જો $\Sigma$ મોડલ તંત્ર હોય, તો યથાર્થતાનો પ્રમેય સ્થાપિત કરે છે કે
$\Sigma \Proves !A$ હોય તો $\Sigma$નાં બધાં !!{formula}sના બધા
દૃષ્ટાંતો માન્ય હોય તેવા નિદર્શનોના કોઈપણ વર્ગ~$\mClass{C}$માં
$!A$ માન્ય છે. ખાસ કરીને, $\Log{K} \Proves !A$ હોય તો $!A$ બધાં
નિદર્શનોમાં સત્ય છે; $\Log{KT} \Proves !A$ હોય તો $!A$ બધાં
સ્વપ્રતિબિંબિત નિદર્શનોમાં સત્ય છે; $\Log{KD} \Proves !A$ હોય તો
$!A$ બધાં સિરિયલ નિદર્શનોમાં સત્ય છે, વગેરે.

પૂર્ણતા એ યથાર્થતાનું વિપરીત વિધાન છે: \Log{K} પૂર્ણ છે એટલે, દાખલા
તરીકે, !!a{formula}~$!A$ માન્ય હોય તો $\Proves !A$. પૂર્ણતા
સાબિત કરવી યથાર્થતા સાબિત કરવા કરતાં ઘણી વધુ અઘરી છે. પહેલાં તેનું
પ્રતિપક્ષી વિધાન ધ્યાનમાં લેવું ઉપયોગી છે: \Log{K} પૂર્ણ છે જો અને
માત્ર જો $\Proves/ !A$ હોય ત્યારે વિરોધી નિદર્શન, એટલે કે એવું
નિદર્શન~$\mModel{M}$ હોય કે $\mSat/{M}{!A}$. સમકક્ષ રીતે ($!A$નું
નિષેધન કરીને), $\Proves/ \lnot !A$ હોય ત્યારે $!A$નું નિદર્શન
હોય છે એવું આપણે સાબિત કરી શકીએ. આવું નિદર્શન રચતી વખતે આપણે
$!A$માં રહેલી માહિતી વાપરી શકીએ છીએ. ચોક્કસ !!{formula}s માટે
નિદર્શન શોધતી વખતે પણ આપણે ઘણી વાર એવું જ કરીએ છીએ: દાખલા તરીકે,
$p \lif \Box q$નું વિરોધી નિદર્શન શોધવું હોય તો તેમાં એવું જગત
હોવું જોઈએ જ્યાં $p$ સત્ય હોય અને $\Box q$ અસત્ય હોય. વળી જ્યાં
$\Box q$ અસત્ય હોય ત્યાંથી પ્રાપ્ય એવું જગત હોવું જોઈએ જ્યાં $q$
અસત્ય હોય. આપણને એટલું જ જાણવું જરૂરી છે: કયા જગતોમાં કયા
!!{propositional variable}s સત્ય છે અને કયા જગતોમાંથી કયા જગતો
પ્રાપ્ય છે.

જોકે પૂર્ણતા સાબિત કરતી વખતે નિદર્શન રચવા માટે આપણી પાસે કોઈ
ચોક્કસ !!{formula}~$!A$ નથી. આપણે સ્થાપિત કરવું છે કે
$\Proves/[\Sigma] \lnot !A$ હોય તેવા દરેક $!A$નું નિદર્શન હોય.
આ લઘુતમ આવશ્યકતા છે, કારણ કે \emph{જો} $\Proves[\Sigma] \lnot !A$
હોય, તો યથાર્થતા મુજબ $!A$નું ($\Sigma$ સત્ય હોય તેવું) કોઈ
નિદર્શન નથી. હવે નોંધો કે $\Proves/[\Sigma] \lnot !A$ જો અને માત્ર
જો $!A$ એ $\Sigma$-સુસંગત હોય. (યાદ કરો કે $\Sigma
\Proves/[\Sigma] \lnot !A$ અને $!A \Proves/[\Sigma] \lfalse$
સમકક્ષ છે.) તેથી આપણું કાર્ય દરેક $\Sigma$-સુસંગત !!{formula}
માટે નિદર્શન રચવાનું છે.

આપણે એ યુક્તિ વાપરીશું કે $!A$ ધરાવતો એવો $\Sigma$-સુસંગત
!!{formula}sનો ગણ શોધવો જેમાં $!A$ને સત્ય બનાવતું જગત કેવું હોવું
જોઈએ તે જણાવતાં અન્ય સૂત્રો પણ હોય. આવા ગણો \emph{પૂર્ણ}
$\Sigma$-સુસંગત ગણો છે. $!A$ને સત્ય બનાવવા એક જ જગતવાળું નિદર્શન
રચવું પૂરતું નથી; તેમાં અનેક જગતો અને પ્રાપ્યતાનો સંબંધ હોવો પડશે.
$!A$ ધરાવતા પૂર્ણ $\Sigma$-સુસંગત ગણમાં $\Box !B$ અને
$\Diamond !C$ સ્વરૂપનાં અન્ય !!{formula}s પણ હશે. બધાં પ્રાપ્ય
જગતોમાં $!B$ સત્ય હોવું જોઈએ અને ઓછામાં ઓછા એકમાં $!C$ સત્ય
હોવું જોઈએ. આ માટે આપણે જગતોના ગણના આધાર તરીકે \emph{બધા}
શક્ય પૂર્ણ $\Sigma$-સુસંગત ગણો લઈશું. કઠિન ભાગ એ નક્કી કરવાનો
છે કે આપણા નિદર્શનમાં એક પૂર્ણ $\Sigma$-સુસંગત ગણ બીજામાંથી
ક્યારે પ્રાપ્ય ગણવો.

આ રીતે વ્યાખ્યાયિત નિદર્શનમાં $!A$ એક જગત પર સત્ય છે જો અને માત્ર
જો $!A$ તે ગણનો !!a{element} છે—એ જગત પોતે પણ એક પૂર્ણ
$\Sigma$-સુસંગત ગણ છે—એવું આપણે બતાવીશું. $!A$ એ
$\Sigma$-સુસંગત હોય તો તે ઓછામાં ઓછા એક પૂર્ણ $\Sigma$-સુસંગત
ગણનો !!a{element} હશે (આ હકીકત આપણે સાબિત કરીશું); તેથી એવું
જગત હશે જ્યાં $!A$ સત્ય છે. આમ આપણને એવું એક જ નિદર્શન મળશે જેમાં
દરેક $\Sigma$-સુસંગત !!{formula}~$!A$ કોઈક જગત પર સત્ય છે.
આ એક જ નિદર્શન $\Sigma$ માટેનું \emph{કૅનોનિકલ} નિદર્શન છે.

\end{document}
